% =====================================================================
\chapter{Muestreo y estimadores}
\label{tema:2}
% =====================================================================

La estadística inferencial trata los datos como valores observados de variables aleatorias cuya distribución depende de parámetros desconocidos. Este tema introduce el modelo de muestra aleatoria simple, la noción de estimador y los criterios para juzgar su calidad: sesgo, varianza, error estándar, error cuadrático medio y consistencia. Se demuestran los dos resultados centrales del muestreo: la media muestral es un estimador insesgado de la media y la varianza muestral con denominador $n-1$ lo es de la varianza. Se sigue el capítulo~6 de Wasserman~\cite{wasserman}, «Models, Statistical Inference and Learning», la §6.3 de Blitzstein y Hwang~\cite{bh}, «Sample moments», la clase~12 del curso Stats~200 de Stanford~\cite{stats200} y la lectura~10a del curso 18.05 del MIT~\cite{mit1805}.

\section{Recordatorio: esperanza y varianza}

Se recogen las propiedades de la esperanza y la varianza de variables aleatorias (discretas o continuas) que se usarán continuamente~\cite[cap.~4 y §7.3]{bh}. Todas las variables que aparecen se suponen con varianza finita.

\begin{proposicion}[Propiedades de la esperanza y la varianza]\label{prop:t2-props}
Sean $X,Y,X_1,\dots,X_n$ variables aleatorias y $a,b\in\R$.
\begin{enumerate}
  \item Linealidad: $\E[aX+bY]=a\E[X]+b\E[Y]$ y, en general, $\E\bigl[\sum_iX_i\bigr]=\sum_i\E[X_i]$.
  \item $\Var(X)=\E[(X-\E X)^{2}]=\E[X^{2}]-(\E X)^{2}$, y $\Var(aX+b)=a^{2}\Var(X)$.
  \item $\Var\bigl(\sum_iX_i\bigr)=\sum_i\Var(X_i)+2\sum_{i<j}\Cov(X_i,X_j)$.
  \item Si $X$ e $Y$ son independientes, $\E[XY]=\E[X]\E[Y]$ y, por tanto, $\Cov(X,Y)=0$. En consecuencia, si $X_1,\dots,X_n$ son independientes, $\Var\bigl(\sum_iX_i\bigr)=\sum_i\Var(X_i)$.
  \item La covarianza es bilineal,
  \[
  \Cov\Bigl(\sum_ia_iX_i,\sum_jb_jY_j\Bigr)=\sum_{i,j}a_ib_j\Cov(X_i,Y_j),
  \]
  y $\Cov(X,X)=\Var(X)$.
\end{enumerate}
\end{proposicion}

\begin{proof}
(i) y la independencia en (iv) son propiedades de la esperanza que se suponen conocidas~\cite[§4.2 y §7.3]{bh}. (ii) Desarrollando el cuadrado y usando (i): $\E[(X-\mu)^{2}]=\E[X^{2}]-2\mu\E[X]+\mu^{2}=\E[X^{2}]-\mu^{2}$ con $\mu=\E X$; además $aX+b-\E[aX+b]=a(X-\mu)$. (v) Se sigue de la linealidad de la esperanza aplicada a $\bigl(\sum_ia_i(X_i-\E X_i)\bigr)\bigl(\sum_jb_j(Y_j-\E Y_j)\bigr)$. (iii) Es (v) con $Y_j=X_j$ y todos los coeficientes iguales a~1, separando los términos $i=j$ de los $i\neq j$. La última afirmación de (iv) se sigue de (iii).
\end{proof}

\section{Muestra aleatoria, parámetro y estadístico}

\begin{definicion}[Muestra aleatoria simple]\label{def:t2-mas}
Una \emph{muestra aleatoria simple} de tamaño $n$ de una distribución $F$ es una colección de variables aleatorias $X_1,\dots,X_n$ independientes y todas con distribución $F$; se dice que son \emph{independientes e idénticamente distribuidas} (\iid) y se escribe $X_1,\dots,X_n\sim F$ \iid\ Los datos $x_1,\dots,x_n$ son una realización (valores observados) de la muestra.
\end{definicion}

El modelo describe el muestreo con reemplazamiento de una población, o el muestreo sin reemplazamiento cuando la población es mucho mayor que la muestra, y también la repetición independiente de un experimento: lanzar una moneda, generar una respuesta de un modelo de lenguaje con una semilla distinta o elegir al azar una pregunta de un banco muy grande.

\begin{definicion}[Modelo paramétrico, parámetro y estadístico]\label{def:t2-parametro}
Un \emph{modelo paramétrico} es una familia de distribuciones $\{F_\theta:\theta\in\Theta\}$ indexada por un \emph{parámetro} $\theta$ (un número o un vector) que toma valores en el \emph{espacio paramétrico} $\Theta$. Un \emph{estadístico} es cualquier función $T=g(X_1,\dots,X_n)$ de la muestra que no depende de parámetros desconocidos; es, por tanto, una variable aleatoria, y su distribución se llama \emph{distribución muestral} de $T$.
\end{definicion}

Por ejemplo, $\{\Ber(p):p\in[0,1]\}$, $\{\Pois(\lambda):\lambda>0\}$ y $\{\Nor(\mu,\sigma^{2}):\mu\in\R,\sigma>0\}$ son modelos paramétricos. Los estadísticos más importantes son la \emph{media muestral} y la \emph{varianza muestral}
\[
\Xbar=\frac1n\sum_{i=1}^{n}X_i,\qquad S^{2}=\frac{1}{n-1}\sum_{i=1}^{n}(X_i-\Xbar)^{2},
\]
que son las cantidades del \cref{tema:1} vistas como variables aleatorias. Se reserva la mayúscula para la variable aleatoria y la minúscula ($\xbar$, $s^{2}$) para su valor observado.

\section{Estimadores y sus propiedades}

\begin{definicion}[Estimador, sesgo, error estándar y error cuadrático medio]\label{def:t2-estimador}
Un \emph{estimador} de un parámetro $\theta$ es un estadístico $\widehat\theta=\widehat\theta(X_1,\dots,X_n)$ que se usa para aproximar $\theta$; su valor con unos datos concretos es una \emph{estimación}. Se definen:
\begin{enumerate}
  \item el \emph{sesgo}: $\sesgo(\widehat\theta)=\E_\theta[\widehat\theta]-\theta$; el estimador es \emph{insesgado} si su sesgo es nulo para todo $\theta\in\Theta$;
  \item el \emph{error estándar}: $\ee(\widehat\theta)=\sqrt{\Var_\theta(\widehat\theta)}$;
  \item el \emph{error cuadrático medio}: $\ECM(\widehat\theta)=\E_\theta\bigl[(\widehat\theta-\theta)^{2}\bigr]$.
\end{enumerate}
El subíndice $\theta$ indica que las esperanzas se calculan con la distribución $F_\theta$.
\end{definicion}

El sesgo mide el error sistemático y el error estándar, la dispersión del estimador alrededor de su propia media. El error estándar suele depender de $\theta$, que es desconocido; al sustituir $\theta$ por su estimación se obtiene el \emph{error estándar estimado}, que se escribe $\widehat{\ee}$.

\begin{proposicion}[Descomposición sesgo--varianza]\label{prop:t2-ecm}
$\ECM(\widehat\theta)=\sesgo(\widehat\theta)^{2}+\Var(\widehat\theta)$.
\end{proposicion}

\begin{proof}
Sea $m=\E[\widehat\theta]$. Se escribe $\widehat\theta-\theta=(\widehat\theta-m)+(m-\theta)$, donde el segundo sumando es una constante. Entonces
\[
\E[(\widehat\theta-\theta)^{2}]=\E[(\widehat\theta-m)^{2}]+2(m-\theta)\E[\widehat\theta-m]+(m-\theta)^{2}=\Var(\widehat\theta)+0+\sesgo(\widehat\theta)^{2},
\]
porque $\E[\widehat\theta-m]=0$.
\end{proof}

Un estimador ligeramente sesgado puede tener menor $\ECM$ que uno insesgado si su varianza es bastante menor (\cref{eje:t2-ecm-normal}); este compromiso entre sesgo y varianza reaparece en el aprendizaje automático (módulo~8).

\section{La media muestral}

\begin{teorema}[Media y varianza de la media muestral]\label{teo:t2-media}
Sean $X_1,\dots,X_n$ \iid\ con media $\mu$ y varianza $\sigma^{2}$. Entonces
\[
\E[\Xbar]=\mu,\qquad \Var(\Xbar)=\frac{\sigma^{2}}{n},\qquad \ee(\Xbar)=\frac{\sigma}{\sqrt n}.
\]
En particular, $\Xbar$ es un estimador insesgado de $\mu$.
\end{teorema}

\begin{proof}
Por la linealidad de la esperanza, $\E[\Xbar]=\frac1n\sum_i\E[X_i]=\frac1n\cdot n\mu=\mu$. Por la \cref{prop:t2-props}(ii) y (iv), usando la independencia,
\[
\Var(\Xbar)=\frac{1}{n^{2}}\Var\Bigl(\sum_{i=1}^{n}X_i\Bigr)=\frac{1}{n^{2}}\sum_{i=1}^{n}\Var(X_i)=\frac{n\sigma^{2}}{n^{2}}=\frac{\sigma^{2}}{n}.\qedhere
\]
\end{proof}

El error estándar de la media decrece como $1/\sqrt n$: para dividirlo por dos hay que multiplicar el tamaño de la muestra por cuatro. Esta ley de la raíz cuadrada gobierna el coste de casi toda estimación por muestreo, desde las encuestas hasta la evaluación de modelos.

\begin{ejemplo}[Proporción muestral]\label{ej:t2-proporcion}
Si $X_1,\dots,X_n\sim\Ber(p)$ \iid\ (por ejemplo, $X_i=1$ si un modelo responde bien a la pregunta $i$), entonces $\mu=p$ y $\sigma^{2}=p(1-p)$. La \emph{proporción muestral} $\widehat p=\Xbar$ es la fracción de aciertos; es insesgada y $\ee(\widehat p)=\sqrt{p(1-p)/n}$, que se estima por $\widehat{\ee}=\sqrt{\widehat p(1-\widehat p)/n}$. Con $n=1000$ preguntas y $\widehat p=0{,}8$, $\widehat{\ee}=\sqrt{0{,}16/1000}\approx0{,}0126$, es decir, algo más de un punto porcentual.
\end{ejemplo}

\section{La varianza muestral}

\begin{teorema}[Insesgadez de la varianza muestral]\label{teo:t2-s2}
Sean $X_1,\dots,X_n$ \iid\ con media $\mu$ y varianza $\sigma^{2}$, con $n\ge2$. Entonces $\E[S^{2}]=\sigma^{2}$.
\end{teorema}

\begin{proof}
\emph{Paso 1.} Por la \cref{prop:t1-extremal}(ii) con $a=\mu$, aplicada a los valores $X_1,\dots,X_n$,
\[
\sum_{i=1}^{n}(X_i-\Xbar)^{2}=\sum_{i=1}^{n}(X_i-\mu)^{2}-n(\Xbar-\mu)^{2}.
\]
\emph{Paso 2.} Se toman esperanzas. Como $\E[(X_i-\mu)^{2}]=\Var(X_i)=\sigma^{2}$ y, por el \cref{teo:t2-media}, $\E[(\Xbar-\mu)^{2}]=\Var(\Xbar)=\sigma^{2}/n$,
\[
\E\Bigl[\sum_{i=1}^{n}(X_i-\Xbar)^{2}\Bigr]=n\sigma^{2}-n\cdot\frac{\sigma^{2}}{n}=(n-1)\sigma^{2}.
\]
\emph{Paso 3.} Dividiendo por $n-1$ se obtiene $\E[S^{2}]=\sigma^{2}$.
\end{proof}

La demostración explica el $n-1$: las desviaciones se miden respecto de $\Xbar$, que está más cerca de los datos que la verdadera media $\mu$ (es la que minimiza la suma de cuadrados), de modo que $\sum_i(X_i-\Xbar)^{2}$ es en promedio menor que $\sum_i(X_i-\mu)^{2}$, exactamente en una cantidad $\sigma^{2}$.

\begin{corolario}[Sesgo de la varianza con denominador $n$]\label{cor:t2-sesgo}
El estimador $\widehat\sigma^{2}=\frac1n\sum_i(X_i-\Xbar)^{2}=\frac{n-1}{n}S^{2}$ cumple $\E[\widehat\sigma^{2}]=\frac{n-1}{n}\sigma^{2}$, luego tiene sesgo $-\sigma^{2}/n$: subestima la varianza, aunque el sesgo tiende a cero cuando $n\to\infty$.
\end{corolario}

\begin{observacion}
Que $S^{2}$ sea insesgado para $\sigma^{2}$ no implica que $S$ lo sea para $\sigma$: en general $\E[S]<\sigma$ (\cref{eje:t2-s-sesgo}). La insesgadez no se conserva por transformaciones no lineales.
\end{observacion}

\section{Consistencia y ley de los grandes números}

Un estimador razonable debe acercarse al parámetro cuando crece el tamaño de la muestra. Se escribe $\widehat\theta_n$ para subrayar la dependencia de $n$.

\begin{definicion}[Consistencia]\label{def:t2-consistencia}
Una sucesión de variables aleatorias $Y_n$ \emph{converge en probabilidad} a una constante $c$, y se escribe $Y_n\convp c$, si para todo $\varepsilon>0$
\[
\lim_{n\to\infty}\PP\bigl(\abs{Y_n-c}>\varepsilon\bigr)=0.
\]
Un estimador $\widehat\theta_n$ es \emph{consistente} si $\widehat\theta_n\convp\theta$ para todo $\theta\in\Theta$.
\end{definicion}

\begin{lema}[Desigualdades de Markov y de Chebyshev]\label{lem:t2-chebyshev}
Si $Y\ge0$ y $a>0$, entonces $\PP(Y\ge a)\le\E[Y]/a$. En consecuencia, si $W$ tiene media $\mu_W$ y varianza finita, para todo $\varepsilon>0$
\[
\PP(\abs{W-\mu_W}\ge\varepsilon)\le\frac{\Var(W)}{\varepsilon^{2}}.
\]
\end{lema}

\begin{proof}
Sea $\ind_A$ la variable que vale $1$ si $Y\ge a$ y $0$ en otro caso. Como $Y\ge0$, se tiene $Y\ge a\ind_A$ en todos los casos, y tomando esperanzas $\E[Y]\ge a\PP(Y\ge a)$. Para la segunda desigualdad se aplica la primera a $Y=(W-\mu_W)^{2}$ y $a=\varepsilon^{2}$, observando que $\abs{W-\mu_W}\ge\varepsilon$ equivale a $(W-\mu_W)^{2}\ge\varepsilon^{2}$.
\end{proof}

\begin{proposicion}[Un $\ECM$ que tiende a cero implica consistencia]\label{prop:t2-ecm-consistencia}
Si $\ECM(\widehat\theta_n)\to0$ para todo $\theta$, entonces $\widehat\theta_n$ es consistente. En particular, basta que el sesgo y la varianza tiendan a cero.
\end{proposicion}

\begin{proof}
Por la desigualdad de Markov aplicada a $Y=(\widehat\theta_n-\theta)^{2}$ con $a=\varepsilon^{2}$, $\PP(\abs{\widehat\theta_n-\theta}>\varepsilon)\le\PP\bigl((\widehat\theta_n-\theta)^{2}\ge\varepsilon^{2}\bigr)\le\ECM(\widehat\theta_n)/\varepsilon^{2}\to0$. La segunda afirmación se sigue de la \cref{prop:t2-ecm}.
\end{proof}

\begin{teorema}[Ley débil de los grandes números]\label{teo:t2-lgn}
Si $X_1,X_2,\dots$ son \iid\ con media $\mu$ y varianza finita $\sigma^{2}$, entonces $\Xbar_n\convp\mu$; es decir, la media muestral es un estimador consistente de la media.
\end{teorema}

\begin{proof}
$\Xbar_n$ es insesgado y $\ECM(\Xbar_n)=\Var(\Xbar_n)=\sigma^{2}/n\to0$; se aplica la \cref{prop:t2-ecm-consistencia}. Explícitamente, $\PP(\abs{\Xbar_n-\mu}>\varepsilon)\le\sigma^{2}/(n\varepsilon^{2})$.
\end{proof}

La cota de Chebyshev es válida para cualquier distribución, pero es muy conservadora; en el \cref{tema:3} el teorema central del límite dará aproximaciones mucho más precisas.

\begin{observacion}[Errores frecuentes]\label{obs:t2-errores}
(a) Confundir el parámetro $\mu$ (un número fijo y desconocido) con el estadístico $\Xbar$ (una variable aleatoria) o con su valor observado $\xbar$ (un número conocido).
(b) Confundir la desviación típica de la población, $\sigma$, con el error estándar de la media, $\sigma/\sqrt n$.
(c) Pensar que «insesgado» significa «cercano al parámetro»: un estimador insesgado con mucha varianza puede ser casi siempre malo.
(d) Usar $\Var(X+Y)=\Var X+\Var Y$ sin comprobar la independencia (o, al menos, la incorrelación).
\end{observacion}

\begin{observacion}[Conexión con el curso]
El módulo~5 retoma estos conceptos (estimadores, sesgo y error cuadrático medio, ley de los grandes números y método de Monte Carlo) y el módulo~8 desarrolla la descomposición sesgo--varianza del error de predicción. En los módulos~22 y~23, la exactitud de un modelo en un benchmark es una proporción muestral (\cref{ej:t2-proporcion}) y la métrica pass@$k$ se estima con un estimador insesgado que se construye con un argumento de muestreo (\cref{eje:t2-passk}).
\end{observacion}

\begin{lecturas}
\item Wasserman~\cite{wasserman}, cap.~6 «Models, Statistical Inference and Learning» (estimación puntual, sesgo, error estándar, consistencia) y cap.~5 «Convergence of Random Variables».
\item Blitzstein y Hwang~\cite{bh}, §6.3 «Sample moments» y §10.2 «Law of large numbers».
\item Stats~200~\cite{stats200}, clase~12 «Parametric models, method of moments estimation» (sesgo, varianza y error cuadrático medio).
\item Orloff y Bloom~\cite{mit1805}, lectura~10a «Introduction to Statistics».
\end{lecturas}

\section*{Ejercicios del tema 2}
\addcontentsline{toc}{section}{Ejercicios del tema 2}

\begin{ejercicio}{1}{eje:t2-conceptos}
Se quiere conocer la exactitud real $p$ de un modelo de lenguaje en un banco de un millón de preguntas. Se eligen al azar $500$ preguntas, el modelo acierta $410$ y se informa de que «la exactitud es del $82\,\%$». Identifique la población, la muestra, el parámetro, el estimador y la estimación. ¿Es aleatorio el número $0{,}82$? ¿Y $\widehat p$?
\end{ejercicio}

\begin{ejercicio}{2}{eje:t2-enumeracion}
Una población está formada por los valores $1,3,5$, cada uno con probabilidad $1/3$. Se extraen muestras de tamaño $n=2$ con reemplazamiento. (a) Calcule $\mu$ y $\sigma^{2}$. (b) Enumere las nueve muestras posibles y obtenga la distribución muestral de $\Xbar$; compruebe que $\E[\Xbar]=\mu$ y $\Var(\Xbar)=\sigma^{2}/2$. (c) Obtenga la distribución de $S^{2}$ y compruebe que $\E[S^{2}]=\sigma^{2}$. ¿Cuánto vale $\E[\widehat\sigma^{2}]$ con denominador $n$?
\end{ejercicio}

\begin{ejercicio}{2}{eje:t2-tres-estimadores}
Sean $X_1,X_2,X_3$ \iid\ con media $\mu$ y varianza $\sigma^{2}$. Se consideran $T_1=\Xbar$, $T_2=(X_1+2X_2+X_3)/4$, $T_3=X_1+X_2-X_3$ y $T_4=(X_1+X_2+X_3)/4$. (a) Calcule el sesgo, la varianza y el $\ECM$ de cada uno. (b) ¿Cuál es el mejor de los insesgados? (c) ¿Para qué valores de $\mu$ tiene $T_4$ menor $\ECM$ que $T_1$?
\end{ejercicio}

\begin{ejercicio}{2}{eje:t2-poisson}
Sean $X_1,\dots,X_n\sim\Pois(\lambda)$ \iid, y $\widehat\lambda=\Xbar$. (a) Demuestre que $\widehat\lambda$ es insesgado y calcule su varianza, su error estándar y su $\ECM$. (b) En $n=50$ bloques de texto se cuentan los errores ortográficos y se obtiene $\xbar=3{,}2$. Dé la estimación de $\lambda$ y su error estándar estimado.
\fuente{adaptado de Stats~200, Stanford~\cite{stats200}, clase~12, ejemplo~12.5.}
\end{ejercicio}

\begin{ejercicio}{2}{eje:t2-expo-sesgo}
Sean $X_1,\dots,X_n\sim\Expo(\lambda)$ \iid\ ($n\ge2$), con media $1/\lambda$. Se admite que $\E[1/\Xbar]=\dfrac{n\lambda}{n-1}$. (a) Calcule el sesgo de $\widehat\lambda=1/\Xbar$. (b) Proponga un estimador insesgado de $\lambda$ de la forma $c/\Xbar$. (c) Con $n=5$ y $\sum x_i=10$, calcule ambas estimaciones.
\fuente{adaptado de Stats~200, Stanford~\cite{stats200}, clase~12, ejemplo~12.6.}
\end{ejercicio}

\begin{ejercicio}{2}{eje:t2-sumas-solapadas}
Sean $X_1,X_2,\dots$ \iid\ con distribución exponencial de tasa $2$ (varianza $1/4$) y sea $n\ge8$. Se definen $X=X_1+\dots+X_n$ e $Y=X_{n-7}+X_{n-6}+\dots+X_{2n-8}$. Calcule $\Cov(X,Y)$ y $\Cor(X,Y)$.
\fuente{adaptado de Orloff y Bloom, MIT 18.05~\cite{mit1805}, \emph{problem set}~5, problema~5.}
\end{ejercicio}

\begin{ejercicio}{2}{eje:t2-benchmark}
Un modelo acierta $300$ de $400$ preguntas elegidas al azar de un banco grande. (a) Estime la exactitud y su error estándar. (b) ¿Cuántas preguntas harían falta para reducir el error estándar a la mitad, si la exactitud estimada no cambia?
\end{ejercicio}

\begin{ejercicio}{2}{eje:t2-combinar}
Se dispone de dos estimadores independientes e insesgados $A$ y $B$ de un mismo parámetro, con $\Var(A)=4$ y $\Var(B)=1$. (a) Demuestre que $T_w=wA+(1-w)B$ es insesgado para todo $w\in\R$. (b) Halle el $w$ que minimiza $\Var(T_w)$ y el valor mínimo. (c) Generalice a varianzas $\sigma_A^{2}$ y $\sigma_B^{2}$.
\end{ejercicio}

\begin{ejercicio}{2}{eje:t2-chebyshev-n}
Se quiere estimar la exactitud $p$ de un modelo con una muestra de $n$ preguntas de modo que $\PP(\abs{\widehat p-p}\ge0{,}05)\le0{,}05$ sea cual sea $p$. Usando la desigualdad de Chebyshev y la cota $p(1-p)\le1/4$, halle un $n$ suficiente.
\end{ejercicio}

\begin{ejercicio}{3}{eje:t2-ecm-normal}
Para una muestra \iid\ de una normal se sabe que $\Var(S^{2})=2\sigma^{4}/(n-1)$. Calcule el $\ECM$ de $S^{2}$ y el de $\widehat\sigma^{2}=\frac{n-1}{n}S^{2}$, y demuestre que el estimador sesgado $\widehat\sigma^{2}$ tiene menor $\ECM$ para todo $n\ge2$.
\end{ejercicio}

\begin{ejercicio}{3}{eje:t2-uniforme}
Sean $X_1,\dots,X_n\sim\Unif(0,\theta)$ \iid\ y $M=\max_iX_i$. (a) Demuestre que $T_1=2\Xbar$ es insesgado y que $\Var(T_1)=\theta^{2}/(3n)$. (b) Demuestre que $\PP(M\le x)=(x/\theta)^{n}$ para $0\le x\le\theta$, calcule $\E[M]$ y $\E[M^{2}]$ y deduzca que $T_2=\frac{n+1}{n}M$ es insesgado con $\Var(T_2)=\dfrac{\theta^{2}}{n(n+2)}$. (c) Demuestre que $\Var(T_2)\le\Var(T_1)$, con igualdad solo si $n=1$.
\fuente{adaptado de Stats~200, Stanford~\cite{stats200}, \emph{homework}~5, problema~3, y de Orloff y Bloom, MIT 18.05~\cite{mit1805}, \emph{problem set}~6, problema~3.}
\end{ejercicio}

\begin{ejercicio}{3}{eje:t2-max-consistente}
Con la notación del \cref{eje:t2-uniforme}, demuestre que $M$ es un estimador consistente de $\theta$ de dos maneras: (a) calculando directamente $\PP(\abs{M-\theta}>\varepsilon)$; (b) mediante la \cref{prop:t2-ecm-consistencia}.
\end{ejercicio}

\begin{ejercicio}{3}{eje:t2-cov-desv}
Sean $X_1,\dots,X_n$ \iid\ con varianza $\sigma^{2}$. Demuestre que $\Cov(\Xbar,X_i-\Xbar)=0$ para cada $i$. (Para datos normales, este hecho es la clave de la independencia entre $\Xbar$ y $S^{2}$.)
\end{ejercicio}

\begin{ejercicio}{3}{eje:t2-s-sesgo}
Demuestre que $\E[S]\le\sigma$, con igualdad solo si $S$ es constante con probabilidad~1. (\emph{Indicación:} $\Var(S)\ge0$.)
\end{ejercicio}

\begin{ejercicio}{3}{eje:t2-passk}
Un modelo genera $n$ soluciones independientes a un problema, cada una correcta con probabilidad $p$, y sea $C\sim\Bin(n,p)$ el número de correctas. Para $1\le k\le n$, la métrica pass@$k$ es la probabilidad de que al menos una de $k$ soluciones independientes sea correcta, $1-(1-p)^{k}$. Demuestre que
\[
\widehat{\pi}_k=1-\binom{n-C}{k}\Big/\binom{n}{k}
\]
es un estimador insesgado de pass@$k$. (\emph{Indicación:} interprete $\binom{n-C}{k}/\binom nk$ como la probabilidad de que un subconjunto de $k$ soluciones elegido al azar no contenga ninguna correcta.)
\end{ejercicio}
